\section{The exact PPR--RPPR comparison bridge}
\label{sec:comparison-bridge}

The common operator makes a comparison possible, but only after accounting
for regularization bias and optimization error.  This section gives that
accounting from first principles.

Recall
\begin{equation}
 f(x)=\frac12x^TQx-b^Tx,
 \qquad
 F_\rho(x)=f(x)+\alpha\rho\|D^{1/2}x\|_1,
 \qquad
 x^\rho:=\argmin_x F_\rho(x).
 \label{eq:ppr-rppr-objectives}
\end{equation}
The unregularized solution is $x^0=Q^{-1}b$.

\begin{lemma}[Sign and residual box for RPPR]
\label{lem:rppr-residual-box}
For every $\rho\geq0$, the unique minimizer satisfies $x^\rho\geq0$.  Its
linear-system residual
\begin{equation}
 r^\rho:=b-Qx^\rho
 \label{eq:rppr-linear-residual}
\end{equation}
obeys the componentwise box
\begin{equation}
 0\leq r^\rho\leq\alpha\rho D^{1/2}\mathbf1.
 \label{eq:rppr-residual-box}
\end{equation}
On every coordinate in $\operatorname{supp}(x^\rho)$, the upper inequality
is an equality.
\end{lemma}

\begin{proof}
The matrix $Q$ is SPD with nonpositive off-diagonal entries and $b\geq0$.
Replacing any $x$ by $|x|$ does not increase the diagonal quadratic or the
$\ell_1$ term, weakly decreases every off-diagonal quadratic term, and does
not increase $-b^Tx$.  Uniqueness therefore gives $x^\rho\geq0$.

The KKT equation is
\begin{equation}
 Qx^\rho-b+\alpha\rho D^{1/2}z=0,
 \qquad
 z_i=1\ \hbox{if }x_i^\rho>0,
 \quad z_i\in[-1,1]\ \hbox{if }x_i^\rho=0.
 \label{eq:rppr-kkt-subgradient}
\end{equation}
If $x_i^\rho=0$, then
$(Qx^\rho)_i=\sum_{j\ne i}Q_{ij}x_j^\rho\leq0\leq b_i$, so
$r_i^\rho\geq0$.  If $x_i^\rho>0$, the KKT equation gives
$r_i^\rho=\alpha\rho\sqrt{d_i}>0$.  Equation
\eqref{eq:rppr-kkt-subgradient} now bounds every coordinate above by
$\alpha\rho\sqrt{d_i}$ and proves the equality statement.
\end{proof}

\begin{theorem}[Exact regularization-bias bridge]
\label{thm:rppr-ppr-bias}
For every $\rho\geq0$,
\begin{equation}
 0\leq x^0-x^\rho\leq\rho D^{1/2}\mathbf1,
 \qquad
 \|x^0-x^\rho\|_{D,\infty}\leq\rho.
 \label{eq:rppr-ppr-bias}
\end{equation}
Consequently, any approximation $\widehat x$ satisfying
$\|\widehat x-x^\rho\|_{D,\infty}\leq\tau$ is a semantic PPR approximation
with
\begin{equation}
 \|\widehat x-x^0\|_{D,\infty}\leq\rho+\tau.
 \label{eq:rppr-ppr-triangle}
\end{equation}
\end{theorem}

\begin{proof}
Because $Q$ is a nonsingular $M$-matrix, $Q^{-1}\geq0$ entrywise.  The shared
normalized Laplacian annihilates $D^{1/2}\mathbf1$, hence
\begin{equation}
 QD^{1/2}\mathbf1=\alpha D^{1/2}\mathbf1,
 \qquad
 Q^{-1}D^{1/2}\mathbf1=\frac1\alpha D^{1/2}\mathbf1.
 \label{eq:degree-ground-state}
\end{equation}
Using \Cref{lem:rppr-residual-box},
\begin{equation*}
 x^0-x^\rho=Q^{-1}r^\rho
 \in\left[0,
 Q^{-1}\alpha\rho D^{1/2}\mathbf1\right]
 =\left[0,\rho D^{1/2}\mathbf1\right].
\end{equation*}
This proves \eqref{eq:rppr-ppr-bias}; the triangle inequality proves
\eqref{eq:rppr-ppr-triangle}.
\end{proof}

This theorem is why an RPPR solver can be used for PPR: choose, for example,
$\rho=\tau=\varepsilon_{\rm ppr}/2$.  It is also why one may not silently
write $\rho=\varepsilon_{\rm ppr}$ inside an RPPR complexity theorem and
ignore inexact-solve error.

\begin{corollary}[Objective-gap conversion]
\label{cor:objective-gap-conversion}
If
\begin{equation}
 F_\rho(\widehat x)-F_\rho(x^\rho)\leq\delta,
 \label{eq:rppr-objective-gap}
\end{equation}
then
\begin{equation}
 \|\widehat x-x^0\|_{D,\infty}
 \leq\rho+\sqrt{\frac{2\delta}{\alpha}}.
 \label{eq:objective-to-semantic}
\end{equation}
In particular, $\rho\leq\varepsilon_{\rm ppr}/2$ and
$\delta\leq\alpha\varepsilon_{\rm ppr}^2/8$ suffice for the semantic PPR
target.
\end{corollary}

\begin{proof}
The objective $F_\rho$ is $\alpha$-strongly convex because
$Q\succeq\alpha I$.  Therefore
\begin{equation*}
 \frac\alpha2\|\widehat x-x^\rho\|_2^2
 \leq F_\rho(\widehat x)-F_\rho(x^\rho)\leq\delta.
\end{equation*}
All graph degrees are at least one, so
$\|z\|_{D,\infty}\leq\|z\|_2$.  Combine this fact with
\Cref{thm:rppr-ppr-bias}.
\end{proof}

The conversion is deliberately conservative: it uses only strong convexity
and not graph-specific Green-function cancellation.  It is nevertheless a
valid common accuracy contract under which two RPPR algorithms can be ranked.

\subsection{On a known face, the \texorpdfstring{$\ell_1$}{l1} term becomes a shifted load}

Let $A=\operatorname{supp}(x^\rho)$.  Positivity and the active KKT equations
give
\begin{equation}
 Q_{AA}x_A^\rho
 =b_A-\alpha\rho D_A^{1/2}\mathbf1_A.
 \label{eq:rppr-active-face-system}
\end{equation}
The inactive KKT test is
\begin{equation}
 0\leq b_j-Q_{jA}x_A^\rho
 \leq\alpha\rho\sqrt{d_j},
 \qquad j\notin A.
 \label{eq:rppr-inactive-face-test}
\end{equation}
Thus, after the correct positive face has been identified, RPPR is not a
different numerical linear-algebra problem: it is the same SPD Stieltjes
principal matrix $Q_{AA}$ with the right-hand side shifted by the linear
$\ell_1$ term.  SOR, CG, or a direct response method can solve
\eqref{eq:rppr-active-face-system}.  What the fixed-face solve does not supply
is the set $A$ or the exterior inequalities \eqref{eq:rppr-inactive-face-test}.

This is the precise answer to the concern ``SOR is non-$\ell_1$, but FISTA is
$\ell_1$-regularized.''  SOR is an iterative method for a linear system, not
an objective class.  It can solve the active linear equations of either PPR
or RPPR.  The nonsmoothness matters at face discovery and at zero-coordinate
KKT tests, not inside a fixed positive face.

\subsection{Three legitimate experimental comparisons}

There are three clean ways to compare the algorithms.

\begin{enumerate}[leftmargin=*]
\item \textbf{Same unregularized problem.}  Set $\rho=0$ so the proximal map
      in FISTA becomes the identity, run both methods on the same graph and
      seed, and stop both at the same semantic PPR error.  The published RPPR
      leaf-star obstruction no longer applies automatically.
\item \textbf{Same regularized problem.}  Fix the same $\rho$, require the
      same objective gap or semantic error via
      \eqref{eq:objective-to-semantic}, and charge both algorithms for finding
      and checking their active coordinates.  A face-aware SOR method must pay
      for \eqref{eq:rppr-inactive-face-test}.
\item \textbf{Same supplied face.}  Give both methods the exact $A$ and
      compare only their restricted numerical work.  This is a linear-solver
      rate experiment, not a locality experiment.
\end{enumerate}

The current center-SOR star and leaf-FISTA star do none of these three
matching operations.  Their correct joint use is diagnostic: one example
proves that signed local acceleration can work, and the other identifies a
specific momentum-locality hazard.
